\model{Queue and Deque}
  \begin{center}
    \small
    \begin{tabular}{p{3.1in} p{2.7in}}
      \begin{minipage}{3.1in}
        \begin{cpplst}
#include <queue>

queue<string> waitlist;
waitlist.push("Alice");
waitlist.push("Bob");
waitlist.push("Carol");

while (!waitlist.empty()) {
    cout << waitlist.front() << endl;
    waitlist.pop();
}
        \end{cpplst}
      \end{minipage}
      &
      \begin{minipage}{2.7in}
        \vskip 6pt
        {\bf waitlist queue (FIFO):}\\[4pt]
        $\rightarrow$\begin{tabular}{|c|c|c|}
          \hline
          Alice & Bob & Carol \\
          \hline
        \end{tabular}$\leftarrow$\\
        \hspace*{0.1in}{\small front (served first)} \hfill {\small back (added here)}\hspace*{0.3in}\\[8pt]
        {\bf Output:}
        \hrule\vskip 5pt
        {\tt Alice}\\
        {\tt Bob}\\
        {\tt Carol}
      \end{minipage}
    \end{tabular}
    \vskip 14pt
    \begin{tabular}{p{3.1in} p{2.7in}}
      \begin{minipage}{3.1in}
        \begin{cpplst}
#include <deque>

deque<int> dq;
dq.push_back(3);
dq.push_front(2);
dq.push_front(1);
dq.push_back(4);

for (int val : dq) {
    cout << val << " ";
}
cout << endl;
dq.pop_front();
dq.pop_back();
        \end{cpplst}
      \end{minipage}
      &
      \begin{minipage}{2.7in}
        \vskip 6pt
        {\bf dq after all push operations:}\\[4pt]
        \begin{tabular}{|c|c|c|c|}
          \hline
          1 & 2 & 3 & 4 \\
          \hline
        \end{tabular}\\[8pt]
        {\bf Output:}
        \hrule\vskip 5pt
        {\tt 1 2 3 4}\\[8pt]
        {\bf dq after pop\_front and pop\_back:}\\[4pt]
        \begin{tabular}{|c|c|}
          \hline
          2 & 3 \\
          \hline
        \end{tabular}
      \end{minipage}
    \end{tabular}
    \vskip 14pt
    \renewcommand{\arraystretch}{1.3}
    \begin{tabular}{|l|c|c|c|c|c|c|}
      \hline
      \rowcolor{orange!20}
      \textbf{Container} & \textbf{push\_front} & \textbf{push\_back} & \textbf{pop\_front} & \textbf{pop\_back} & \textbf{at(i)} & \textbf{range-for} \\
      \hline
      \tt vector & no  & yes & no  & yes & yes & yes \\
      \hline
      \tt list   & yes & yes & yes & yes & no  & yes \\
      \hline
      \tt queue  & no  & yes & yes & no  & no  & no  \\
      \hline
      \tt deque  & yes & yes & yes & yes & yes & yes \\
      \hline
    \end{tabular}
  \end{center}

  {\it\large Refer to Model 4 above as your team develops consensus answers
    to the questions below.}

  \quest{15 min}

  \Q A {\tt queue} models a waiting line: items are added at the back
    and removed from the front.  This behavior is called
    {\it FIFO} (First In, First Out).
    \begin{enumerate}
      \itemsep 10pt
      \item Who was the first person pushed onto the {\tt waitlist}?
        \hfill\ans[1.5in]{"Alice"}

      \item Who was printed (served) first?
        \hfill\ans[1.5in]{"Alice"}

      \item If you called {\tt waitlist.push("Dave")} before the
        {\tt while} loop, what would be the last name printed?
        \hfill\ans[1.5in]{"Dave"}

      \item Why can't you use a range-based for loop to print the
        contents of a {\tt queue}?  (Hint: look at the table.)
        \begin{answer}[0.5in]
          A \cpp{queue} does not support range-based for iteration.
          It only exposes the front element; elements must be removed
          one at a time with \cpp{pop()}.
        \end{answer}
    \end{enumerate}

  \vskip -20pt

  \Q Using the operations table, compare {\tt queue} and {\tt deque}.\key\\[-2mm]
    \begin{enumerate}
      \itemsep 10pt
      \item What does the ``d'' in {\tt deque} stand for?
        \hfill\ans[2in]{``double-ended''}

      \item List two operations that {\tt deque} supports but {\tt queue}
        does not.
        \begin{answer}[0.5in]
          \cpp{push_front}, \cpp{pop_back}, \cpp{at(i)}, and
          range-based for.  (Any two are acceptable.)
        \end{answer}

      \item A {\tt deque} supports {\tt at(i)} but a {\tt queue} does not.
        What does this tell you about which container supports random access?
        \begin{answer}[0.5in]
          Only \cpp{deque} supports random access by index.
          A \cpp{queue} only allows access to the front element.
        \end{answer}
    \end{enumerate}

  \vskip -20pt

  \Q For each scenario below, choose the most appropriate container
    from {\tt vector}, {\tt list}, {\tt map}, {\tt set}, {\tt queue},
    or {\tt deque}.  Justify each choice.
    \begin{enumerate}
      \itemsep 12pt
      \item Storing a list of student exam scores that you need to
        access by index and sort.
        \begin{answer}[0.5in]
          \cpp{vector} --- supports random access by index and works
          well with sorting algorithms.
        \end{answer}

      \item Keeping track of customers waiting in line at a support desk
        (serve in the order they arrived).
        \begin{answer}[0.5in]
          \cpp{queue} --- FIFO behavior naturally models a waiting line.
        \end{answer}

      \item Storing a dictionary mapping English words to their
        definitions.
        \begin{answer}[0.5in]
          \cpp{map} --- stores key-value pairs and supports fast lookup
          by key (word).
        \end{answer}

      \item Recording which usernames have already been taken (no
        duplicates, fast lookup).
        \begin{answer}[0.5in]
          \cpp{set} --- stores unique values and supports fast
          membership testing with \cpp{count}.
        \end{answer}

      \item Simulating a deck of cards where cards can be drawn or
        added from either end.
        \begin{answer}[0.5in]
          \cpp{deque} --- supports efficient push and pop from both
          ends.
        \end{answer}
    \end{enumerate}

  \vskip -10pt

  \Q The file {\tt activity28d.cpp} contains the queue and deque
    examples.  Run it and verify the output matches the model.
    Then write C++ code using a {\tt deque<int>} that adds the
    values 1, 2, 3 to the front (in order) and 4, 5, 6 to the
    back, then prints all values.
    \begin{answer}[1.25in]
      \begin{cpplst}
deque<int> dq;
dq.push_front(3);
dq.push_front(2);
dq.push_front(1);
dq.push_back(4);
dq.push_back(5);
dq.push_back(6);
for (int val : dq) {
    cout << val << " ";
}
cout << endl;
      \end{cpplst}
    \end{answer}
